\section{Introduction}

\subsection{Motivation}

\paragraph{Challenges of power grid stability require better power system control:} Modern power grids need a robust Model-Based Systems Engineering (MBSE) method to connect control design with real-world control deployment and routine tuning. Recent incidents, including the inter-area oscillations observed in Europe in 2016 and 2017 \cite{ENTSOE_2016, ENTSOE_2017}, highlight how crucial proper tuning of a Power System Stabilizer (PSS) is to preserve adequate damping throughout the system. More recently, the Iberian grid incident of 2025 \cite{ENTSOE_Iberian_Factual_2025, ENTSOE_Iberian_2026} revealed complex 0.63 Hz and 0.2 Hz oscillations. Although initiated by inverter-based resources, the incident report emphasizes that poor coordination and the absence of detailed models for active damping controls—including existing PSS—limited operators’ ability to foresee and contain the disturbance, forcing rapid revisions to operating rules.

Discrepancies between simulations and field behavior often stem from outdated modeling workflows that fail to capture actual hardware logic. Such control-induced issues have been increasingly observed in inverter-based resources, including irregular dynamics in wind plants due to setpoint truncation \cite{Pudasaini2026} and non-sinusoidal responses in solar PV installations \cite{Mishra2026}. This fragility extends to traditional synchronous generation; for example, a 1.5 Hz forced oscillation event was recently triggered during a reactive power capability test at a conventional power plant in Dominion \cite{Mishra2024_Greensville}. Although the power plant model was accurate, the PSS was not designed for the specific operating condition encountered, leading to sustained oscillations. These challenges demand adaptable controllers that can be redesigned and deployed throughout their operational lifecycle \cite{Bombois2024, Lauria2024_SEGAN}.

\paragraph{Development and deployment challenges of power system controllers hinder adoption of novel stabilizers:} When developing new control strategies to address grid stability issues, the path from new control designs to Controller Hardware-in-the-Loop (CHiL) tests is currently inefficient due to a lack of dynamic model portability and broken traceability. Offline power system simulation tools, such as PSS\regtm{}E or PSCAD\regtm{}, cannot be directly used for real-time simulation. This necessitates manual re-implementation, which creates ambiguity and consumes significant time and resources. Furthermore, control code is often hardware-dependent; for example, specialized wide-area PSS implementations frequently require bespoke code for specific platforms such as National Instruments\regtm{} hardware \cite{Rebello2020}. This fragmented workflow and the high cost of real-time simulators and prototyping hardware, such as CompactRIO\regtm{} (National Instruments), create significant barriers to deploying novel control strategies.

\paragraph{eFMI tooling for Modelica provides new opportunities for power system design and control:} The acausal, equation-based Modelica\regtm{}\footnote{\label{fnt:modelica}\tiny{\url{https://modelica.org/language/}}} modeling language has been successfully used to design and simulate power systems and grids~\cite{ipsl,openipsl}. Its object-oriented, system-level design is particularly well-suited for rapid prototyping of design variants and modifications~\cite{Burger2019}. Modelica also provides the means of modeling clocked subsystems, such as digital controllers, seamlessly integrated within the continuous simulation of plant models \cite{Elmqvist2012}. This makes Modelica an excellent choice for physics-based modeling and simulation of power systems with novel stabilizing controllers. Recent support for the new eFMI\regtm{} Standard\footnote{\label{fnt:efmi-website}\tiny{\url{https://www.efmi-standard.org/}}} \cite{eFMI_Standard} by Modelica tools (such as Dymola\regtm{}\footnote{\label{fnt:dymola}\tiny{\url{https://www.3ds.com/products/catia/dymola}}}) extends the applicability of physics models beyond offline simulation to the embedded domain. The eFMI capabilities provided by Dymola facilitate the deployment of Modelica-based models in safety-critical applications that operate under stringent hard real-time constraints.
\subsection{Contributions}

This paper pioneers a Modelica and eFMI-based approach to design power system damping controllers. We use Modelica and eFMI to (1) rapidly prototype and conduct Model- and Software-in-the-Loop (MiL, SiL) testing of PSS designs, (2) deploy the designed models (plant and controller models) directly onto microcontrollers, and finally (3) verify in CHiL tests that embedded behavior aligns with the offline simulations. Our approach solves many of the previously bespoke design, development, and deployment challenges associated with a PSS:

\squishlist
    \item Common physics-based models are the single source of truth, reducing model-misalignment and inconsistency errors, and preventing behavior that violates or neglects laws of physics.
    \item The MBSE approach reduces errors during the transition from physics-based models to actual embedded solutions, as is typical for manually developed control code or code based on heterogeneous tooling at different abstraction levels that requires manual adaptation within deep technology stacks.
    \item Directly deploying models (controller and plant) on microcontrollers enables rapid control prototyping and reduces the risk of late detection of system-integration issues. Expensive real-time platforms can be avoided and the complexity of the tooling stacks involved can be reduced.
    \item Gaps in model portability and traceability are minimized.
\squishend

\noindent{}To validate our approach, we designed real-time-suited Modelica models based on the Open-Instance Power System Library\footnote{\label{fnt:openipsl}\tiny{\url{https://github.com/OpenIPSL/OpenIPSL}}} (\texttt{OpenIPSL}), generated the corresponding eFMI solutions in Dymola, system-integrated them on STM32\regtm{} microcontrollers with the help of STM32CubeMX\footnote{\label{fnt:stm32cubemx}\tiny{\url{https://www.st.com/en/development-tools/stm32cubemx.html}}} and STM32CubeIDE\footnote{\label{fnt:stm32cubeide}\tiny{\url{https://www.st.com/en/development-tools/stm32cubeide.html}}}, and systematically validated each step in respective MiL, SiL, and CHiL tests. All materials are provided open-source under the 3-Clause BSD License, as an extensible workbench framework\footnote{\label{fnt:pss-repository}\tiny{\url{https://github.com/ALSETLab/efmi-pss-use-case}}} ultimately empowering researchers to rapidly prototype grid controllers with guaranteed traceability and significantly reduced capital expenditure. Our approach provides a low-cost proof-of-concept platform for PSS microcontroller prototyping.

\subsection{Paper structure}

The remainder of this paper is organized as follows: Section \ref{sec:efmi-and-use-case-workflow-overview} provides an overview of the eFMI Standard and the proposed MBSE workflow. Section \ref{sec:modeling-for-mil-testing} details the modeling approach and MiL validation using the new \texttt{OpenIPSL\_CHIL} library. Section \ref{sec:efmuGenAndSiL} discusses the synthesis of eFMUs and SiL tests and code analyses for such, while Section \ref{sec:STM32-integration-of-eFMUs} describes the system integration on STM32 hardware. Real-time simulation results and CHiL performance analyses are presented in Section \ref{sec:rtchil}. Section \ref{sec:conclu} provides conclusions and outlines future work, followed by an Appendix summarizing the related work that highlights the limitations of current practices in the power industry related to damping controls.
